\documentclass[10pt,twocolumn,letterpaper]{article}
\usepackage[letterpaper,left=0.62in,right=0.62in,top=0.68in,bottom=0.72in]{geometry}
\usepackage{amsmath,booktabs,graphicx,microtype,times,url}
\graphicspath{{../outputs/proposal_models_5m/}{../outputs/armas_5m_full_study/}}
\setlength{\columnsep}{0.23in}
\title{Physics-Guided Temporal Transformer for Aviation\\Effective Dose-Rate Estimation}
\author{Sanjib K C\\Department of Physics and Astronomy, Georgia State University\\
\texttt{[add email]} \quad PantherID: \texttt{[add PantherID]}}
\date{}
\begin{document}
\maketitle

\begin{abstract}
This preliminary study executes the model family specified in the project
proposal. The 5-minute ARMAS data contain 92,415 observations in three
flight-separated partitions, 35 environmental variables over 12 time steps,
and eight current flight/geomagnetic variables. P1 fits preprocessing and
neural weights, P2 is the validation set, and P3 is the final test set. This
completed preliminary assignment is denoted 123.
Six PyTorch experiments compare a 43-input joint multilayer perceptron (JMLP),
a CNN--MLP, a direct Transformer, and physics-guided residual Transformers.
P2 validation selects the JMLP overall; on P3 it reaches RMSE 3.778 versus
3.990 $\mu$Sv/h for NAIRAS V3. A residual Transformer trained with high-dose
weight four gives the best RMSE above 15 $\mu$Sv/h, 5.650 versus 6.503 for
NAIRAS. Thus current inputs remain a strong selection choice, while
imbalance-aware residual attention benefits the safety-relevant upper tail.
\end{abstract}

\section{Introduction}
ARMAS supplies in-flight measurements of aviation effective dose rate, whereas
NAIRAS V3 supplies a physics-based estimate \cite{tobiska2016,mertens2013}.
Prior static learning reduced error, especially when NAIRAS underpredicts large
ARMAS observations \cite{sadykov2026,kc2025}. This study asks whether the
preceding 55 minutes add information, which variables and lags drive a
correction, and whether high-dose weighting improves accuracy above
15 $\mu$Sv/h without unacceptable degradation elsewhere.

The physics-guided target is
\begin{equation}
r_i=A_i-N_i,\qquad \hat A_i=N_i+f_\theta(X_i,N_i),
\end{equation}
where $A_i$ is ARMAS and $N_i$ is NAIRAS V3. The model corrects a physical
estimate rather than learning dose rate entirely from scratch.

\section{Related Work}
Tobiska et al. introduced ARMAS for global, real-time aviation radiation
monitoring \cite{tobiska2016}. Mertens et al. developed and validated NAIRAS as
a physics-based aircraft dose model \cite{mertens2013}. The machine-learning-ready
dataset aligns ARMAS and NAIRAS with aircraft position, neutron-monitor data,
solar-wind variables, particle fluxes, and geomagnetic indices
\cite{sadykov2026}.

Earlier static machine learning showed that geospace properties can improve
aviation-radiation nowcasts relative to uncorrected NAIRAS \cite{kc2025}. This
project extends that work by testing the preceding hour with self-attention,
following the Transformer framework of Vaswani et al. \cite{vaswani2017}.

\section{Model and Method}
\subsection{Data and Leakage Control}
Each record contains a $12\times35$ temporal tensor and eight static variables.
Source arrays place the current value at index zero; each array is reversed once
to chronological order from $t-55$ to $t$ minutes. Vehicle ID is excluded.
Median imputation and standardization use P1 statistics only. The preliminary
assignment is 123, where the digits give the train--validation--test order. All
P1 records supply optimization, P2 controls early stopping and model selection,
and P3 is scored only after checkpoint selection is complete.

\begin{table}[t]
\centering
\caption{Partition sizes for P1 train, P2 validation, and P3 test.}
\begin{tabular}{lrrr}\toprule
Partition & Rows & $A_i>15$ & Percent \\ \midrule
P1 & 32,369 & 5,324 & 16.4 \\
P2 & 30,752 & 4,983 & 16.2 \\
P3 & 29,294 & 4,947 & 16.9 \\ \bottomrule
\end{tabular}
\end{table}

\begin{figure}[t]
\centering\includegraphics[width=\columnwidth]{figure01_target_distribution.png}
\caption{ARMAS target distribution using the full P1+P2+P3 dataset.}
\end{figure}

\subsection{Deep-Learning Models}
The JMLP concatenates eight static and 35 current environmental inputs, then
uses 128- and 64-unit GELU layers with normalization and dropout. The hybrid
CNN--MLP applies two one-dimensional convolutions over the history and joins
the pooled temporal representation to a current-input branch.

The Transformer projects each 35-variable step to 48 dimensions, adds
sinusoidal positional encoding and a learned summary token, and applies two
four-head pre-normalized encoder blocks. A separate MLP embeds current inputs.
Direct models predict standardized ARMAS; residual models also receive NAIRAS
and predict standardized $A_i-N_i$.

All networks use AdamW, cosine learning-rate decay, Huber loss, gradient
clipping at 1.0, batch size 512, at most 24 epochs, and patience five.
High-dose residual variants multiply losses above 15 $\mu$Sv/h by two or four.

\begin{table}[t]
\centering\small
\caption{Executed architectures and trainable parameter counts.}
\begin{tabular}{lrr}\toprule
Model & Parameters & Epochs \\ \midrule
JMLP & 14,209 & 11 \\
CNN--MLP & 30,401 & 13 \\
Direct Transformer & 61,345 & 11 \\
Residual Transformer & 61,441 & 8 \\
Residual Transformer $w=2$ & 61,441 & 10 \\
Residual Transformer $w=4$ & 61,441 & 9 \\ \bottomrule
\end{tabular}
\end{table}

\begin{figure}[t]
\centering\includegraphics[width=\columnwidth]{02_learning_curves.png}
\caption{P2 validation loss used for early stopping and checkpoint selection.}
\end{figure}

\section{Experiments and Results}
\subsection{Evaluation Protocol}
We report MAE, RMSE, coefficient of determination, and bias overall and for
$A_i\leq15$ and $A_i>15$ $\mu$Sv/h. The central comparison is NAIRAS V3,
JMLP, CNN--MLP, direct Transformer, residual Transformer, and two
imbalance-aware residual Transformers. Group occlusion replaces one
standardized temporal variable or complete lag with its P1 mean and measures
the P2 RMSE change. Occlusion measures model dependence rather than causality.

\subsection{P3 Test Results}
\begin{table}[t]
\centering\small
\caption{P3 test results. RMSE is in $\mu$Sv/h; lower is better.}
\begin{tabular}{lrrr}\toprule
Model & All & $A_i>15$ & $R^2$ \\ \midrule
NAIRAS V3 & 3.990 & 6.503 & .449 \\
JMLP & 3.778 & 6.571 & .506 \\
CNN--MLP & 4.239 & 7.076 & .378 \\
Direct Transformer & 3.889 & 6.874 & .477 \\
Residual Transformer & 3.811 & 6.602 & .497 \\
Residual Transformer $w=2$ & \textbf{3.712} & 6.086 & \textbf{.523} \\
Residual Transformer $w=4$ & 3.822 & \textbf{5.650} & .495 \\ \bottomrule
\end{tabular}
\end{table}

\begin{figure}[t]
\centering\includegraphics[width=\columnwidth]{01_proposal_model_rmse.png}
\caption{Overall and high-dose P3 test RMSE. Ranking changes when the upper tail is emphasized.}
\end{figure}

P2 validation selects the JMLP overall. On P3 it reaches RMSE 3.778, MAE 2.793,
and $R^2=0.506$, a 5.3\% RMSE reduction from NAIRAS. The weight-two residual
Transformer has the lowest descriptive P3 RMSE (3.712), but this test-set
ranking is not used to revise selection. The CNN--MLP is weaker than NAIRAS.

The upper tail supports the imbalance-aware hypothesis. Increasing the
residual loss weight from one to four reduces high-dose RMSE from 6.602 to
5.650 and high-dose bias from $-5.064$ to $-3.513$ $\mu$Sv/h. Overall RMSE
increases only from 3.811 to 3.822. The operational objective therefore changes
which model should be selected.

\begin{figure}[t]
\centering\includegraphics[width=\columnwidth]{03_best_observed_vs_predicted.png}
\caption{Observed versus predicted P3 dose rate for the model selected on P2.}
\end{figure}

\subsection{Temporal Interpretation}
Feature-group occlusion gives the largest RMSE increases for the Newark
neutron monitor (0.030), Oulu neutron monitor (0.025), Thule neutron monitor
(0.011), and solar F10.7 flux (0.009 $\mu$Sv/h). No lag occlusion increases
RMSE; the least negative change is at the current observation. Small and negative lag effects indicate strong
temporal redundancy and help explain why attention does not dominate the JMLP.

\begin{figure}[t]
\centering\includegraphics[width=\columnwidth]{04_transformer_occlusion.png}
\caption{Residual-Transformer variable and lag occlusion. Positive values indicate worse RMSE after removal.}
\end{figure}

\section{Conclusions}
The study supports learned correction of a physics model, but temporal
complexity is not automatically useful. The 32,369-row training partition may
also be small for a 61,000-parameter attention model, while closely spaced
environmental variables are highly redundant. The weighted residual
Transformer is valuable when high-dose underprediction is the stated objective.

These are single-seed, single-split preliminary results. Measurements within
flights may be correlated, P3 includes 4,947 high-dose cases, and occlusion can
be distorted by correlated inputs. CPU constraints motivated compact
architectures. P3 results must not be used for further tuning.

\section{Future Work}
The completed preliminary analysis reports assignment 123: P1 train, P2
validation, and P3 test. The future robustness study will use the five remaining
train--validation--test permutations: 132, 213, 231, 312, and 321. Every
rotation will fit preprocessing and model weights only on its training
partition, select configurations on its validation partition, and evaluate its
test partition only after selection. JMLP, direct Transformer, residual
Transformer, and weighted residual Transformer will be repeated with at least
three seeds. Integrated gradients and attention diagnostics will complement
grouped occlusion, with interpretations reported as associations.

\section{Reproducibility}
The executed script \texttt{armas\_proposal\_models\_5m.py} saves checkpoints,
per-epoch histories, predictions, metrics, plots, parameter counts, occlusion
tables, and a JSON run manifest. Random seed 42 and the PyTorch version are
recorded. Repository: \url{https://github.com/thesanjibkc/Deep_learning_project}.

\section{Peer Evaluation}
This is an individual project; therefore, there are no team peers to evaluate.

\clearpage
\begin{thebibliography}{9}
\bibitem{tobiska2016} W. K. Tobiska et al., ``Global real-time dose measurements using the Automated Radiation Measurements for Aerospace Safety system,'' \emph{Space Weather}, 14, 1053--1080, 2016.
\bibitem{mertens2013} C. J. Mertens et al., ``NAIRAS aircraft radiation model development, dose climatology, and initial validation,'' \emph{Space Weather}, 11, 603--635, 2013.
\bibitem{sadykov2026} V. M. Sadykov et al., ``Machine Learning-Ready Data Sets for the Analysis and Nowcasting of Atmospheric Radiation at Aviation Altitudes,'' arXiv:2602.07108, 2026.
\bibitem{kc2025} S. K C, V. M. Sadykov, and D. Kempton, ``Nowcasting of Aviation Radiation Using Geospace Environment Properties: A Machine Learning Approach,'' arXiv:2511.15746, 2025.
\bibitem{vaswani2017} A. Vaswani et al., ``Attention Is All You Need,'' \emph{NeurIPS}, 2017.
\end{thebibliography}
\end{document}
